% fig_appendix_fig7_oracle - generated by bin/make_figures.py (bin/make_figures.py sha256:054c165eac74)
% Plan label: [ORACULO: as-submitted (O1): umbral y detector elegidos con etiquetas de evaluacion] [EXPLORATORIO] (audit made before the analysis plan was fixed)
% numbers_used.md ids (1): R.FIG7.AGG.BARS_LE_TRIVIAL
% 216 values read from fig7_valores_publicados.csv (216 = 9 x 12 x 2).
\begin{figure*}[t]
  \centering
  \includegraphics[width=\textwidth]{figures/fig_appendix_fig7_oracle.pdf}
  \caption{Oracle results of the submitted regime, redrawn from the 216 bar values of its figure (nine datasets, twelve variants, two metrics); not comparable with the main results. For every dataset and variant the threshold was chosen on the whole stream with the evaluation labels, and the detector shown is the one with the lower WindowDiff, again chosen with those labels: an asterisk on the axis label marks the Centroid Detector (KMD), and the other bars show the Distribution Detector (MVD). Grey bars are the F1 and hatched bars the WindowDiff, both against the extended labels of that regime (`Smooth' in the legend), the latter from the submitted, non-standard implementation ($k=\lfloor 2\cdot 10^{10}/N \rfloor$ records, window sums instead of boundary counts, 8-bit overflow; Section~5.6). MLP-AR and MLP-TSR are replicates of one per-record model (Section~4.3.2). The dashed line is the F1 of the always-alarm classifier against the same extended labels, with its value under each panel title, and downward triangles mark the 24 of the 108 F1 bars that do not exceed it, an exploratory count.}
  \label{fig:appendix-fig7}
\end{figure*}
